% =====================================================================
%  Formato de página según la normativa ICV-611 (UCM).
%  Cada bloque indica la regla de 01_normativa/REGLAS_FORMATO.md que aplica.
%  LIN = Lineamientos (§6), TPL = template Word oficial.
% =====================================================================

% --- Márgenes: sup 2,5 / inf 2,5 / izq 3 / der 3 cm [§1, LIN] ---
% El encabezado y el número de página quedan DENTRO del margen (como en TPL).
\geometry{
  top        = 2.5cm,
  bottom     = 2.5cm,
  left       = 3cm,
  right      = 3cm,
  headheight = 15pt,
  headsep    = 0.85cm,
  footskip   = 1.2cm,
}

% --- Interlineado 1,5 [§1, LIN] ---
% "1,5 líneas" de Word con Times 12 = 20,7 pt entre líneas. En LaTeX eso es
% \setstretch{1.43} (el \onehalfspacing de LaTeX es más estrecho, ~18 pt, y
% haría que la tesis ocupe menos páginas que su equivalente en Word).
\setstretch{1.43}

% --- Párrafos: justificado, sangría 1,25 cm, sin espacio extra [§1, TPL] ---
\setlength{\parindent}{1.25cm}
\setlength{\parskip}{0pt}

% --- Numeración y encabezados [§2, LIN + TPL] ---
% Número: abajo a la derecha, solo el número.
% Encabezado: arriba a la derecha "Capítulo N. Nombre"; en anexos "Anexos".
% Sin líneas ni leyendas extra.
\renewcommand{\headrulewidth}{0pt}
\renewcommand{\footrulewidth}{0pt}
\fancypagestyle{cuerpo}{%
  \fancyhf{}%
  \fancyhead[R]{\nouppercase{\leftmark}}%
  \fancyfoot[R]{\thepage}%
}
\fancypagestyle{plain}{%  primera página de cada capítulo y preliminares: solo número
  \fancyhf{}%
  \fancyfoot[R]{\thepage}%
}
\renewcommand{\chaptermark}[1]{\markboth{Capítulo~\thechapter.~#1}{}}
\renewcommand{\sectionmark}[1]{}
\pagestyle{plain}

% --- Títulos [§5, TPL] ---
% Capítulo: "1. Introducción y Planteamiento del Problema", 18 pt negrita, izquierda
\titleformat{\chapter}[hang]
  {\normalfont\fontsize{18}{22}\selectfont\bfseries\raggedright}
  {\thechapter.}{0.5em}{}
% Títulos sin número (Resumen, Índice General, Glosario, Bibliografía...): igual, sin número
\titleformat{name=\chapter, numberless}[hang]
  {\normalfont\fontsize{18}{22}\selectfont\bfseries\raggedright}
  {}{0pt}{}
\titlespacing*{\chapter}{0pt}{0pt}{30pt}
% Sección: "1.1. Título", 14 pt negrita
\titleformat{\section}
  {\normalfont\fontsize{14}{17}\selectfont\bfseries\raggedright}
  {\thesection.}{1em}{}
\titlespacing*{\section}{0pt}{12pt}{3pt}
% Subsección: "1.1.1. Título", 12 pt negrita
\titleformat{\subsection}
  {\normalfont\normalsize\bfseries\raggedright}
  {\thesubsection.}{1em}{}
\titlespacing*{\subsection}{0pt}{12pt}{3pt}
% Máximo 3 niveles numerados (1.1.1) [LIN]
\setcounter{secnumdepth}{2}
\setcounter{tocdepth}{1}       % Índice General: capítulos y secciones (como TPL)

% --- Índices (tocloft) [§3, §5, TPL] ---
\renewcommand{\cftchapfont}{\normalfont}
\renewcommand{\cftchappagefont}{\normalfont}
\renewcommand{\cftchapleader}{\cftdotfill{\cftdotsep}}
\renewcommand{\cftchapaftersnum}{.}
\renewcommand{\cftsecaftersnum}{.}
\setlength{\cftbeforechapskip}{4pt}
\setlength{\cftsecindent}{1.25cm}
\setlength{\cftsecnumwidth}{3em}
% Índice de Figuras / Tablas: "Figura 1.1: Título ....... 11"
\renewcommand{\cftfigpresnum}{Figura~}
\renewcommand{\cftfigaftersnum}{:}
\renewcommand{\cfttabpresnum}{Tabla~}
\renewcommand{\cfttabaftersnum}{:}
\setlength{\cftfigindent}{0pt}
\setlength{\cfttabindent}{0pt}
\setlength{\cftfignumwidth}{6.2em}
\setlength{\cfttabnumwidth}{5.8em}
% Rótulo "Pág." sobre la columna de números (se escribe al inicio de cada índice)
\newcommand{\rotulopag}{\protect\hfill\protect\textbf{Pág.}\protect\par}

% --- Leyendas: 11 pt, negrita solo en el rótulo, "Figura 2.2: Título" [§6, LIN] ---
% Tablas: título ARRIBA. Figuras: título ABAJO.
\DeclareCaptionFont{simple}{\setstretch{1.1}}
\captionsetup{font={small,simple}, labelfont=bf, textfont=normalfont,
  labelsep=colon, justification=centering, skip=6pt}
\captionsetup[table]{position=above}
\captionsetup[figure]{position=below}
% Contenido de las tablas también en 11 pt [§6, LIN]
\AtBeginEnvironment{table}{\small\setstretch{1.15}}
\AtBeginEnvironment{longtable}{\small\setstretch{1.15}}

% --- Numeración por capítulo: ecuaciones (1.1), figuras 1.1, tablas 1.1 [§6] ---
\numberwithin{equation}{chapter}

% --- Anexos: "Anexo I: Título", secciones "I.1.", encabezado "Anexos" [§5, TPL] ---
\newcommand{\iniciaranexos}{%
  \cleardoublepage
  \markboth{Anexos}{Anexos}%
  \chapter*{Anexos}%
  \thispagestyle{cuerpo}%
  \addcontentsline{toc}{chapter}{Anexos}%
  \appendix
  \renewcommand{\thechapter}{\Roman{chapter}}%
  \titleformat{\chapter}[hang]
    {\normalfont\fontsize{18}{22}\selectfont\bfseries\raggedright}
    {\appendixname~\thechapter:}{0.4em}{}%
  \renewcommand{\chaptermark}[1]{\markboth{Anexos}{Anexos}}%
  \assignpagestyle{\chapter}{cuerpo}% en TPL los anexos llevan "Anexos" también en su 1.ª página
  % el formato del índice se escribe DENTRO del .toc para que aplique solo a los anexos
  \addtocontents{toc}{\protect\renewcommand{\protect\cftchappresnum}{\appendixname~}}%
  \addtocontents{toc}{\protect\renewcommand{\protect\cftchapaftersnum}{:}}%
  \addtocontents{toc}{\protect\setlength{\protect\cftchapnumwidth}{5em}}%
}

% --- Bibliografía: título "Bibliografía", en el índice, encabezado "Bibliografía" ---
\defbibheading{bibliografia}{%
  \chapter*{Bibliografía}%
  \addcontentsline{toc}{chapter}{Bibliografía}%
  \markboth{Bibliografía}{Bibliografía}%
}
% Sangría francesa y separación entre referencias (APA)
\setlength{\bibitemsep}{0.5\baselineskip}

% --- Rutas ---
\graphicspath{{figuras/}}

% --- Estilo de código en anexos ---
\lstset{basicstyle=\ttfamily\small, breaklines=true, frame=single, numbers=left, numberstyle=\tiny}
